\documentclass{ximera}
%geschreven 16/11/26
\title{Meten van spanning en spanningsbronnen}
\author{Daan Lipkens}

\begin{document}
\begin{abstract}
    Hoe je een spanning meet met een voltmeter, en welke soorten spanningsbronnen er bestaan.
\end{abstract}
\maketitle
\label{hoofdstuk:appendix_spanning}

\section*{Inleiding}
In het hoofdstuk over potentiële elektrische energie, potentiaal en spanning leerde je dat spanning het verschil in potentiaal is tussen twee punten. In de praktijk meet je een spanning met een voltmeter, en zorgt een spanningsbron ervoor dat er een spanning is. Deze appendix bundelt wat je daarover uit de tweede graad kent:
\begin{itemize}
    \item het meten van spanning met een voltmeter (sectie~\verwijs[Meten van spanning]{sec:spanning_meten}{appendix_spanning_meten});
    \item gelijk- en wisselspanningsbronnen en hun symbolen (sectie~\verwijs[Spanningsbronnen]{sec:spanningsbronnen}{appendix_spanning_bronnen}).
\end{itemize}

\end{document}
